\documentclass[12pt]{exam}
\printanswers	% Pour afficher les solutions
	% On précise ci-dessous l'en-tête d'une question.
	\qformat{\hspace*{0.6cm}\textbf{Question~\thequestion~(\totalpoints~points)}\hfill}
	% On précise ci-dessous l'affichage du numéro de la page.
	\runningfooter{}{}{\scriptsize \sffamily Page \thepage ~de \numpages}
	% Pour que les points des sous-questions soient additionnés à côté de la question.
	\addpoints		
	% Pour que les points des sous-questions apparaissent dans la marge.
	\pointsinmargin
	
\usepackage[T1]{fontenc}		% Permet de tapper directement les accents.
\usepackage[utf8]{inputenc}		% Le type d'encodage UTF-8 est assez commun, d'autres encodages sont possibles.
\usepackage[letterpaper,margin=2cm]{geometry}
\usepackage{enumerate}

\usepackage{hyperref}
\hypersetup{
	colorlinks=true,
	urlcolor=blue, 	% Liens externes
	linkcolor=blue, % Liens internes
	citecolor=red   % Citations
}

\usepackage{microtype}
\usepackage{lmodern}			% Pour changer la police, aucunement requis.
\renewcommand{\rmdefault}{ptm}		% Pour changer la police, aucunement requis.
\usepackage{graphicx}			% Pour inclure des images.
\usepackage[frenchb]{babel}		% Pour signifier à LaTeX que le document est en français.
\usepackage{tikz}			% Pour utiliser tikz, paquet permettant de faire des images.
	\usetikzlibrary{calc}
\usepackage{pgfplots}		% Pour tracer des graphiques avec tikz.
\usepackage{amsmath}		% Pour tapper des math dans LaTeX.
\usepackage{amssymb}		% Pour utiliser certains symboles mathématiques.
\usepackage{multicol}		% Requis pour certains tableaux
\usepackage{relsize}		% Pour utiliser la commande \mathlarger{}

\begin{document}
	\begin{titlepage}
		\begin{center}
			\Huge Devoir 1\\
			\vspace{3cm}
			\Large Réalisé par Personne 1 et Personne 2\\
			\vspace{3cm}
			\large Présenté à Jérôme Soucy\\
			Dans le cadre du cours\\
			MAT-2903: Thèmes mathématiques pour l'enseignement secondaire\\
		\vfill
			\includegraphics[width=45mm]{logo_ul.pdf}\\
			15 octobre 2016
		\end{center}
	\end{titlepage}
	
\begin{questions}
\question Soit une ellipse d'équation
\[\frac{x^2}{a^2}+\frac{y^2}{b^2}=1,\]
et soit $P$ un point lui appartenant. Appelons $T$ la droite tangente à l'ellipse en $P$, et appelons $A$ et $B$ les points d'intersection de la droite $T$ avec l'axe des $x$ et l'axe des $y$ respectivement. Soient maintenant $D_x$ la droite parallèle à l'axe des $x$ passant par $B$, et $D_y$ la droite parallèle à l'axe des $y$ passant par $A$. Finalement, appelons $Q$ le point de rencontre des droites $D_x$ et $D_y$. Remarquons que ces définitions n'ont du sens que si $P$ n'est pas un sommet de l'ellipse. On fera donc cette hypothèse.
  \begin{parts}
	\part[5]\label{q:ggb1} Dans GeoGebra, définissez deux curseurs $a$ et $b$, prenant leurs valeurs dans l'intervalle $(0,5)$, puis construisez l'ellipse d'équation
	\[\frac{x^2}{a^2}+\frac{y^2}{b^2}=1.\]
	\begin{solution}
	 Voir construction sur le \href{https://www.geogebra.org/}{Site GeoGebra}. % Choisir le lien de votre construction.
	\end{solution}

	\part[10]\label{q:ggb2} Étant donné un point $P$ sur l'ellipse, construisez, à l'aide de l'outil {\bf Lieu}, le lieu des points $Q$ décrit dans le préambule de cette question.
	\begin{solution}
	 Voir construction sur le \href{https://www.geogebra.org/}{Site GeoGebra}. % Choisir le lien de votre construction.
	\end{solution}
	\part[5] Montrez que les points de coordonnées $(x,y)$ appartenant au lieu précisé à la sous-question précédente satisfont
	\[\frac{a^2}{x^2}+\frac{b^2}{y^2}=1.\]
Vous pouvez utiliser le fait que l'équation de la tangente à l'ellipse d'équation
\[\frac{x^2}{a^2}+\frac{y^2}{b^2}=1\]
au point $(x_0,y_0)$ est donnée par
\[\frac{xx_0}{a^2}+\frac{yy_0}{b^2}=1.\]
	\begin{solution}
	 Voici la solution.
	\end{solution}
  \end{parts}
  
\question Dans le contexte de l'enseignement secondaire, les coniques étudiées ont généralement leurs axes parallèles à l'axe du repère. Dans le cas contraire, les équations obtenues pour résoudre les différents problèmes associés sont plus complexes. La connaissance d'une formule simple, donnant la distance entre un point et une droite, permet d'obtenir assez facilement l'équation d'une telle conique.
  \begin{parts}
	\part[6] Soit $d$ une droite d'équation $y=mx+b$ dans un repère orthonormée $Oxy$. Soit $P$ un point du plan, extérieur à $d$, de coordonnées $(x_0,y_0)$. Montrez que la distance entre $d$ et $P$ est donnée par
	\[\frac{|mx_0-y_0+b|}{\sqrt{1+m^2}}.\]
	\begin{solution}
	 Voici la solution.
	\end{solution}
	\part[12] Une ellipse a pour foyers les points $F_1$ et $F_2$, de coordonnées respectives $(-6,4)$ et $(-4,2)$. Son excentricité est de $\frac{\sqrt{2}}{2}$. Trouvez l'équation de l'ellipse ainsi que celle de chacune de ses directrices, de même que les coordonnées du centre et de ses quatre sommets.
	\begin{solution}
	 Voici la solution.
	\end{solution}
	\part[2] Construisez l'ellipse dans GeoGebra, en prenant soin d'y tracer les directrices et d'identifier le centre, les foyers, et les quatre sommets. Incluez une image de votre construction dans le devoir.\\
	\begin{solution}
	 Voir construction sur le \href{https://www.geogebra.org/}{Site GeoGebra}. % Choisir le lien de votre construction.
	\end{solution}
  \end{parts}
\end{questions}

\end{document} 